% q0065.tex — To Be Rational or Adaptive?
\question{
  id        = {0065},
  title     = {To Be Rational or Adaptive?},
  source    = {OBECON},
  year      = {2026},
  round     = {Seletiva IEO -- Phase C},
  language  = {English},
  topic     = {Macro},
  subtopic  = {Expectations / Phillips Curve / Monetary Policy},
  type      = {dissertative},
  statement = {%
    Consider an economy with a Central Bank following a monetary rule that
    delivers inflation $\pi_t = \bar{\pi}$ whenever agents hold rational
    expectations $\pi^e_t = \bar{\pi}$.

    The Phillips Curve is:
    \[
      \pi_t - \pi^e_t = -\alpha(u_t - u_n),
    \]
    where $u_t$ is unemployment, $u_n$ is the natural rate, and $\alpha > 0$.

    \begin{enumerate}[(a)]
      \item (30 points) What property must the Central Bank's monetary rule
            satisfy for the rational-expectations equilibrium $\pi^e_t =
            \bar{\pi}$ to hold? Explain the link between credibility and
            expectations.

      \item (30 points) Suppose the government unexpectedly increases
            spending. Describe the effect on inflation and explain what
            interest-rate response the Central Bank should take to stabilise
            inflation back to $\bar{\pi}$.

      \item (40 points) Under what circumstances might agents shift to
            \emph{adaptive} expectations ($\pi^e_t = \pi_{t-1}$)? How does
            this change the effectiveness and cost of monetary policy?
    \end{enumerate}
  },
  choices   = {},
  answer    = {},
  solution  = {%
    \textbf{(a)} The monetary rule must be \textbf{credible}: agents will
    only set $\pi^e_t = \bar{\pi}$ if they are confident the Central Bank
    will deliver on its target. Credibility requires consistent, transparent
    policy backed by institutional independence. Without it, expectations
    become unanchored, reducing policy effectiveness.

    \textbf{(b)} Higher government spending raises aggregate output, reducing
    unemployment below $u_n$. From the Phillips Curve, $u_t < u_n$ implies
    $\pi_t > \pi^e_t = \bar{\pi}$, so inflation rises above target. The
    Central Bank should \textbf{raise interest rates} to reduce investment
    and consumption, increasing unemployment back toward $u_n$ and returning
    inflation to $\bar{\pi}$.

    \textbf{(c)} Agents shift to adaptive expectations when they lose
    confidence in CB commitment --- e.g.\ after repeated failure to hit the
    announced target. Under adaptive expectations, monetary policy becomes
    \textbf{less effective}: agents no longer believe announced targets, so
    disinflationary policy requires prolonged tight monetary policy and is
    more costly in terms of output and unemployment. The sacrifice ratio
    rises, making credibility central to efficient policy transmission.
  },
}
